\bookchapter{Two}{Sequences and Series}

\section{Sequences}

Most of the concepts in analysis can be reduced to statements about the behavior of sequences.  Therefore it is important to have a clear understanding of them before moving forward.




\begin{definition}\textsc{Sequences} are functions whose domain is $\mathbb{N}$.  \end{definition}

The formal definition leads immediately to the familiar depiction of a sequence as an ordered list of real numbers.  Given a function $f: \mathbb{N} \to \mathbb{R}$,   $f(n)$ is just the $n$th term on the list.  The notations for sequences reinforce this understanding.

What is crucial to understand is that a sequence is an \emph{infinite} list of real numbers.  What happens at the beginning is of little importance since the business of analysis is concerned with the behavior of the infinite tail of a given sequence.  

\begin{definition}Convergence of a Sequence:  A sequence $(p_n)$ converges to a real number $p$ if, \marginnote{It might be well to point out that our definition of convergent sequence depends not only on the sequence in question but also on the metric space $X$.}for every positive number $\varepsilon$, there exists an $N \in \mathbb{N}$ such that whenever $n \geq N$ it follows that $|p_n - p|< \varepsilon$. \end{definition}


For Rudin, a sequence $\{p_n\}$ in a metric space $X$ is said to converge if there is a point $p \in X$ with the following property:  For every $\varepsilon > 0$ there is an integer $N$ such that $n \geq N$ implies that $d(p_n, p) < \varepsilon$.

In this case, we also say that  $\{p_n\}$ converges to $p$, or that $p$ is the limit of  $\{p_n\}$ and we write that 
\[\lim_{n \to \infty} p_n = p   \qquad (p_n) \to p\]


If a sequence does not converge, it is said to diverge.  \marginnote{Note: the definition of divergence depends on the negation of quantifiers.  A sequence diverges if, for \emph{every} candidate limit $p$, there exists an $\varepsilon>0$ for which no $N$ works: however large $N$ is, some $n \geq N$ has $d(p_n, p) \geq \varepsilon$.}

\begin{definition}
We recall that the set of all points $p_n  (n=1,2,3, \ldots)$  is the \textsc{range} of the sequence.  The range can be infinite or finite.  


 The sequence $\{p_n\}$ is said to be \textsc{bounded} if its range is bounded.  That is, a sequence is bounded if there exists an $M > 0$ such that $|p_n|\leq M$ for all $n \in \mathbb{N}$.\end{definition}


\bigskip
In an effort to decipher this complicated definition, it helps to consider the phrase ``$d(p_n, p) < \varepsilon$'' and think about the points that satisfy an inequality of this type.


\begin{definition}$\varepsilon$-neighborhoods:  Given a point $p \in X$ and a positive number $\varepsilon > 0$,  the set
\[V_{\varepsilon}(p) = \{x \in X: d(x,p)< \varepsilon\}\]
is called the $\varepsilon$-neighborhood of $p$. \end{definition}

Notice that $V_{\varepsilon}(p)$ consists of all those points whose distance from p is less than $\varepsilon$.  Recasting the definition of convergence in terms of $\varepsilon$-neighborhoods gives us a more geometric impression of what is being described.

\begin{definition}\newthought{Convergence of a Sequence (Topological Version):}  a sequence $\{p_n\}$ in a metric space $X$ converges to $p$, if, given any $\varepsilon$-neighborhood  $V_{\varepsilon}(p)$ of $p$, there exists a point in the sequence after which all the terms are in $V_{\varepsilon}(p)$.  In other words, every $\varepsilon$-neighborhood centered at $p$ contains all but a finite number of the terms of $\{p_n\}$.
\end{definition}

Notice that these two definitions of convergence say exactly the same thing.  The natural number $N$ from before is the point where the sequence enters $V_{\varepsilon}(p)$ and never leaves.  It should be apparent that \emph{the value of $N$ depends on the choice of} $\varepsilon$.   

The smaller the $\varepsilon$-neighborhood, the larger $N$ may have to be.


\noindent \emph{Example:}  We claim that $$ \lim \bigg( \frac{1}{\sqrt{n}} \bigg) = 0 $$

\begin{proof}
Let $\varepsilon >0$ be an arbitrary positive number.  Choose a natural number $N$ such that 

$$	\frac{1}{\sqrt{N}} < \varepsilon, \quad \text{that is,} \quad N > \frac{1}{{\varepsilon}^2}.	$$

We now verify that this choice of $N$ has the desired property.  Let $n \geq N$.  Then,

$$	n> \frac{1}{\varepsilon^2}	  \quad \text{implies} \quad \frac{1}{\sqrt{n}} < \varepsilon \quad \text{and hence} \quad |a_{n} -0|< \varepsilon$$

\end{proof}

\section{Some Examples:}\label{sec:sequence-examples}

\begin{enumerate}
	\item If $s_n = 1/n$ then $\lim_{n \to \infty} s_n = 0$.  The range is infinite and the sequence is bounded.
	\item If $s_n = n^2$ then the sequence is divergent, unbounded and has infinite range.
	\item If $s_n = 1 + [{(-1)^n}/n]$ then the sequence $\{s_n\}$ converges to 1, is bounded, and has infinite range.
	\item If $s_n = i^n$, then the sequence is divergent, is bounded, and has finite range.
	\item If $s_n = 1$ then the sequence converges to 1, is bounded, and has finite range.
	\item  Now take the sequence $a_n = 1/\sqrt{n}$.
	
	Our intuitive understanding tells us that the sequence converges to 0 but let's first explore the relationship between $\varepsilon$ and $N$ in the definition of convergence.
	
	Let's define a sort of ``target zone''  and say that we take $\varepsilon = 1/10$.  By claiming that the limit of the sequence is 0, we are saying that the terms of the sequence eventually get arbitrarily close to 0.  
	
	We have set $\varepsilon = 1/10$ as our standard of ``closeness,'' which leads to the neighborhood $(-1/10, 1/10)$ around the limit 0.  How far out into the sequence before the terms fall into this interval?  The 100th term $a_{100} = 1/10$ puts us right on the boundary so we can say that,  if $n>100$ then $a_n \in (-1/10, 1/10)$.  Thus for $\varepsilon = 1/10$ we choose $N=101$ (or anything larger) as our response.
	
	But our choice of $\varepsilon $ was maybe rather whimsical.  How can we be sure that for every choice of $\varepsilon$ we can find an appropriate $N$?  For this problem, the desired relationship is implicit in the algebra carried out to compute our previous answer.  Whatever $\varepsilon $ happens to be, we want 
	\[\frac{1}{\sqrt{n}} < \varepsilon\]
	which is equivalent to 
	\[n > \frac{1}{\varepsilon ^ 2}\]
	
	Now we are ready to write a formal argument.

We claim that \[\lim_{n \to \infty} \frac{1}{\sqrt{n}} = 0\]

\begin{proof}
Let $\varepsilon> 0 $ be an arbitrary positive number.  Choose a natural number $N$ satisfying $N > \frac{1}{\varepsilon ^ 2}$.  

We now verify that this choice of $N$ has the desired property.  Let $n \geq N$.  Then,
\[n > \frac{1}{\varepsilon^2} \implies \frac{1}{\sqrt{n}}< \varepsilon\]

Then \[\left|\frac{1}{\sqrt{n}} - 0\right| < \varepsilon. \]  \end{proof}
\end{enumerate}



The  definition of convergence given earlier is the result of hundreds of years of refining the intuitive notion of a limit into a mathematically rigorous statement.  The logic of such a statement is intimately tied to the use of the quantifiers ``for all'' and ``there exists.''  The ability to write proofs of convergence depends heavily on a firm understanding of how to manipulate such statements.

\newthought{Template for Proof that $(x_{n}) \to x$}
	\begin{enumerate}
			\item Let $\varepsilon>0$ be arbitrary
			\item  Demonstrate a choice for $N \in \mathbb{N}$.  This usually requires the most work, which should be done before actually writing the proof (less messy)
			\item Now, show that $N$ actually works
				\begin{enumerate}
						\item Assume $n \geq N$
						\item With $N$ well chosen, it should be possible to derive the inequality that $|x_{n} - x| < \varepsilon$.
				\end{enumerate}
	\end{enumerate}


\noindent \emph{Example:}  We claim that $$ \lim \bigg( \frac{n+1}{n} \bigg) = 1 $$

\begin{proof}
Let $\varepsilon >0$ be an arbitrary positive number.  Choose a natural number $N$ such that $ N > \frac{1}{\varepsilon}	$.


We now verify that this choice of $N$ is appropriate.  Let $n \geq N$.  Then,

$$	n> \frac{1}{\varepsilon}	  \quad \text{implies} \quad \frac{1}{n} < \varepsilon \quad \text{and hence} \quad \bigg|  \frac{n+1}{n}  -1  \bigg|  = \bigg|  \frac{1}{n} \bigg|< \varepsilon$$

as desired.
\end{proof}



Significant insight into the role of the quantifiers in the definition of convergence
can be gained by studying an example of a sequence that does not have a limit.

\begin{definition}Divergence:  A sequence that does not converge is said to diverge. \end{definition}


To prove that a particular number $x$ is not the limit of a sequence $(x_{n})$ we must produce a single value of $\varepsilon$ for which no $N \in \mathbb{N}$ works.
\medskip

\noindent
\emph{Exercise 2.2.1.} Verify, using the definition of convergence of a sequence, that
the following sequences converge to the proposed limit.
\begin{enumerate}[(a)]
\item $\lim \frac{1}{(6n^{2}+1)} = 0$.
\item $\lim \frac{(3n+1)}{(2n+5)} = \frac{3}{2}$.
\item $\lim \frac{ 2}{\sqrt{n+3}} = 0$.
\end{enumerate}

\begin{proof}
	\begin{enumerate}[(a)]
			\item Let $\varepsilon > 0$ be arbitrary. We must show that there exists an
$N \in \mathbb{N}$ such that $n\geq N$ implies $\big|\frac{1}{(6n^{2}+1)} - 0 \big|  < \varepsilon$.  Well, manipulation of the statement tells us that 

$$\bigg|\frac{1}{(6n^{2}+1)} \bigg|  = \frac{1}{(6n^{2}+1)}  \quad \text{and therefore that}$$

$$ \frac{1}{(6N^{2}+1)}  < \varepsilon \quad \text{which holds when} \quad N > \sqrt{\frac{1- \varepsilon}{6\varepsilon}}$$

(for $\varepsilon < 1$; any $N$ works when $\varepsilon \geq 1$).  However, we can avoid the square root entirely.  Since $6n^2 + 1 > 6n$, we have $\frac{1}{6n^{2}+1} < \frac{1}{6n}$, so it is enough to choose $N > \frac{1}{6\varepsilon}$: then $n \geq N$ gives $\frac{1}{6n^2+1} < \frac{1}{6n} \leq \frac{1}{6N} < \varepsilon$.  

It follows that $n \geq N$  implies that $\frac{1}{(6n^{2}+1)}  < \varepsilon$ as desired.


\item  Let $\varepsilon>0$.  Produce $N$ such that $n\geq N$ implies
$$   \bigg|\frac{(3n+1)}{(2n+5)} -\frac{3}{2}\bigg| <  \varepsilon $$

A little algebra shows that the left-hand side equals $\frac{13}{4n+10}$, and that $n>\frac{13-10\varepsilon}{4\varepsilon}$ makes this smaller than $\varepsilon$.  So we pick $N >\frac{13-10\varepsilon}{4\varepsilon}$; then for $n \geq N$ it follows that $\frac{13}{4n+10} < \varepsilon$.


  \item Let $\varepsilon>0$.  We must produce an $N$ so that $n \geq N$ implies that $|\frac{2}{\sqrt{n+3}} - 0|< \varepsilon$.  Well,
  $$ \bigg| \frac{2}{\sqrt{n+3}} - 0 \bigg| = \frac{2}{\sqrt{n+3}} 	$$
  
  Pick $N$ such that $N > 4/{\varepsilon^2} - 3$.  It then follows that when $n \geq N$ we get $\frac{2}{\sqrt{n+3}} <\varepsilon$ as desired.
	\end{enumerate}
\end{proof}


\noindent
\emph{Exercise 2.2.4.} Argue that the sequence $$1,0,1,0,0,1,0,0,0,1,0,0,0,0,1,(\text{5 zeros}),1,...$$
does not converge to zero. For what values of $\varepsilon > 0$ does there exist a response $N$. For which values of $\varepsilon > 0$ is there no suitable response?

\begin{proof}
For any $\varepsilon$ that is greater than $1$, there exists a response $N$. In this case, $N$ can be any natural number.

For any $\varepsilon$ that is less than or equal to $1$, there exists no suitable response. This is because, although the 1s in the sequence occur less and less frequently as we go out the sequence, there is still no point in the sequence where the sequence enters the neighborhood $(-\varepsilon, \varepsilon)$ and never leaves.

\end{proof}


\emph{Exercise 2.2.6.} Suppose that for a particular $\varepsilon > 0$ we have found a suitable value of $N$ that ``works'' for a given sequence in the sense of Definition 2.2.
\begin{enumerate}[(a)]	
\item Then, any larger/smaller (pick one) $N$ will also work for the same $\varepsilon > 0$. 
\item Then, this same $N$ will also work for any larger/smaller value of $\varepsilon$.
\end{enumerate}

\begin{proof}
(a) Any larger $N$ will also work for the same $\varepsilon > 0$. (b) This same $N$ will also work for any larger value of $\varepsilon$.
\end{proof}

We now summarize some important properties of convergent sequences in metric spaces.

\begin{theorem}  Let $\{p_n\}$ be a sequence in a metric space $X$.
\begin{enumerate}[(a)]
	\item $\{p_n\}$  converges to $p$ iff every neighborhood of $p$ contains all but finitely many of the terms of $\{p_n\}$
	\item If $p \in X, p' \in X$ and if $\{p_n\}$ converges to $p$ and to $p'$, then $ p=p'$
	\item If $\{p_n\}$ converges, then $\{p_n\}$ is bounded.
	\item If $E \subset X$ and if $p$ is a limit point of $E$, then there is a sequence $\{p_n\}$ in $E$ such that $\{p_n\}$ converges to $p$.
\end{enumerate}
\end{theorem}

\begin{proof}
(a) Suppose that $p_n \to p$ and let $V$ be a neighborhood of $p$.  For some $\varepsilon > 0 $, the conditions $d(p,q)< \varepsilon, \quad q \in X$ imply $q \in V$.  Corresponding to this $\varepsilon$ there exists $N$ such that $n \geq N$ implies $d(p_n, p)< \varepsilon$.  Thus $n \geq N$ implies $p_n \in V$.

Conversely, suppose that every neighborhood of $p$ contains all but finitely many elements of $\{p_n\}$.  Fix $\varepsilon>0$ and let $V$ be the set of all $q \in X$ such that $d(p,q)< \varepsilon$.  By assumption, there exists an $N$ (corresponding to this $V$) such that $p_n \in V$ if $n \geq N$.  Thus $d(p_n, p) < \varepsilon$ if $n \geq N$; hence $p_n \to p$.

(b)  Let $\varepsilon > 0$ be given.  There exist integers $N, N'$ such that
\[n\geq N \implies d(p_n, p) < \frac{\varepsilon}{2}\]
\[n \geq N' \implies d(p_n, p') < \frac{\varepsilon}{2}\]

If we choose $n\geq \max(N, N')$, we have that 
\[d(p,p') \leq d(p_n,p) + d(p_n, p') < \varepsilon\]
Since $\varepsilon$ was an arbitrary choice, we can therefore conclude that $d(p',p) =0$ and that $p=p'$.

(c)  Suppose that $p_n \to p$.  Then there exists an integer $N$ such that $n\geq N$ implies that $d(p_n, p) < 1$.  Now suppose that we set $r$ to
\[r = \max \{1, d(p_1, p), \ldots, d(p_N, p)\}\]

Then $d(p_n, p) \leq r$ for all $n =1, 2, 3, \ldots$.  Therefore the sequence is bounded.

(d) Let $p$ be a limit point of $E \subset X$. \marginnote{Limit point:  A point $p \in X$ is a limit point of the set $E$ if every neighborhood $V_{\varepsilon}(p)$ contains a point $q \neq p$ such that $q \in E$} Then \emph{every} neighborhood of $p$ contains points $p_n\in E$ with $p_n\neq p$. In other words, every open ball around $p$, contains other points of $E$.  

We now build a sequence $\{p_n\}$ in $E$.  We define it as follows:  for each positive integer $n$,  let $p_n$ be a point of $E$, different from $p$, such that $d(p_n, p) < 1/n$.  

Given $\varepsilon > 0$, choose $N$ such that $N > \frac{1}{\varepsilon}$ (which is equivalent to $N\varepsilon>1$).  If $n > N$, it follows that $d(p_n, p)< \frac{1}{n} < \frac{1}{N} < \varepsilon$.  Hence $p_n \to p$.
\end{proof}


\section{The Algebraic and Order Limit Theorems}

The point of having such a logically tight description of convergence is so that we can confidently state and prove statements about convergent sequences in general.  We are ultimately trying to resolve arguments about what is and is not true regarding the behavior of limits with respect to the mathematical manipulations we intend to inflict on them.

\begin{definition}\newthought{Bounded:} A sequence $(x_{n})$ is bounded if there exists a number $M>0$ such that $|x_{n}|\leq M$ for all $n \in \mathbb{N}$. \end{definition}

Geometrically this implies that we can find an interval on the real number line $[-M, M]$ that contains every term of the sequence $(x_{n})$.

\begin{theorem} Every convergent sequence is bounded. \end{theorem}

\begin{proof}
Assume $(x_{n})$ converges to a limit $l$. This means that given a particular value of $\varepsilon$, we know that there exists an $N \in \mathbb{N}$ such that if $n\geq N$, then $x_{n}$ is in the interval $(l-\varepsilon, l+\varepsilon)$.  Without knowing whether or not $l$ is positive or negative, we can nonetheless be sure that \marginnote{ Note that $|x_n| = |x_n + l - l| \leq |x_n - l| + |l| < \varepsilon + |l|$}  $$ |x_{n}| < |l| + \varepsilon $$ for all $n\geq N$.

We still need to worry (just a little) about the terms that come before the Nth term.  Because there are only a finite number of these terms, we can let $$ M = \max \{|x_{1}|,|x_{2}|, \ldots, |x_{N-1}|,|l|+\varepsilon \}$$
It follows that $|x_{n}|\leq M$ for all $n \in \mathbb{N}$ as desired.
\end{proof}

The following theorem  illustrates that sequences behave very well with respect to each other.  This should be easily done using the convergence definition and the properties of norms from earlier.

\begin{theorem}[ Algebraic Limit Theorem] Suppose that $\{s_n\}$ and $\{t_n\}$ are complex sequences, and that they converge to $s$ and $t$ respectively.  Then:
\begin{enumerate}[(a)]
	\item $\{s_n + t_n\} \to s+t$
	\item $\{cs_n\} \to cs$ and $\{c+s_n\} \to c+s$, for any number $c$
	\item $\{s_n t_n\} \to st$
	\item $\{1/s_n\} \to 1/s$, provided that $s_n \neq 0$ for all $n$ and $s \neq 0$
\end{enumerate}
\end{theorem}

\begin{proof}
(a)  Given $\varepsilon > 0$, there exist integers $N_1, N_2$ such that 
\[n \geq N_1 \implies |s_n - s|< \frac{\varepsilon}{2}\]
\[n \geq N_2 \implies |t_n - t|< \frac{\varepsilon}{2}\]

If $N = \max(N_1, N_2)$ then $n\geq N$ implies
\[|(s_n + t_n)-(s + t)|\leq |s_n - s|+|t_n - t|<\varepsilon\] proving that $\{s_n + t_n\} \to s+t$.

(b)  If $c = 0$ there is nothing to prove.  Otherwise, given $\varepsilon>0$, there exists $N$ such that $n\geq N \implies |s_n-s|<\varepsilon/ |c|$.  That is, whenever $n\geq N$

\[|cs_n -cs |= |c||s_n - s|< \varepsilon\]  proving that $\{cs_n\} \to cs$.

Similarly, since $|(c+s_n) - (c+s)| = |s_n - s|$, any $N$ that works for $\{s_n\}$ also shows that $\{c+s_n\} \to c+s$.

(c) Given $\varepsilon > 0$, there exist integers $N_1, N_2$ such that 
\[n \geq N_1 \implies |s_n - s|< \sqrt{\varepsilon}\]
\[n \geq N_2 \implies |t_n - t|< \sqrt{\varepsilon}\]

If we take $N = \max(N_1, N_2), n\geq N$ implies that 
\[|(s_n -s )(t_n -t)| \leq |s_n - s| |t_n -t| < \varepsilon\]

which is to say that the sequence $\{(s_n - s)(t_n - t)\}$ converges to 0.

To show that $\{s_n t_n \} \to st$ we write 


$\begin{array} {lcl} 
|s_n t_n  - st|& = &|(s_n - s) (t_n - t) + s(t_n -t) + t(s_n - s)|\\ 


& \leq & |(s_n - s) (t_n - t)| + |s||(t_n -t) |+ |t||(s_n - s)|    \end{array}$

Therefore for $n  \geq N$ as before, we have that 

\[|(s_n - s) (t_n - t)| + |s||(t_n -t) |+ |t||(s_n - s)| < \varepsilon + |s| \sqrt{\varepsilon} + |t| \sqrt{\varepsilon}.\]
Since $\varepsilon$ was arbitrary, the right-hand side can be made as small as we please, and so $\{s_n t_n\} \to st$.



\emph{Alternative Proof:}  Begin by observing that 


$\begin{array} {lcl} 
|s_n t_n  - st|& = &|s_n t_n - st_n + st_n  -st|\\ 
& \leq & |s_n t_n - st_n| + |st_n  -st|\\
&=	& |t_n||s_n - s| + |s||t_n - t| 
    \end{array}$
    
Letting $\varepsilon$ be arbitrary, we proceed with the strategy of making each piece less than $\varepsilon/2$.  Using the fact that convergent sequences are bounded (proved previously) we know that there exists a bound $M>0$ such that $|t_n|\leq M$ for all $n \in \mathbb{N}$.  Therefore for any $\varepsilon >0$ there exists $N_1, N_2$ such that 
\[n\geq N_1 \implies |s_n - s| <  \frac{1}{M}\frac{\varepsilon}{2}\]
\[n\geq N_2 \implies |t_n - t| <  \frac{1}{|s|}\frac{\varepsilon}{2}\]

(If $s = 0$ the second term vanishes and $N_2$ is not needed.)  Pick $N = \max(N_1, N_2)$ and the calculation shows that 

\smallskip

$\begin{array} {lcl} 
|s_n t_n  - st|& \leq &|t_n||s_n - s| + |s||t_n - t| \\ 
& \leq & M|s_n - s| + |s||t_n - t|\\
& <	&  M\frac{1}{M}\frac{\varepsilon}{2} + |s|\frac{1}{|s|}\frac{\varepsilon}{2}\\
&=& \varepsilon
    \end{array}$

\smallskip



(d) To show that $\{1/s_n\} \to 1/s$ we first observe that 
\[\left\lvert \frac{1}{s_n} - \frac{1}{s} \right\rvert = \frac{|s-s_n|} {|s||s_n|}.\]

Since $\{s_n\} \to s$ we can make the preceding numerator as small as we like by choosing $n$ large.  The problem is that we also need an estimate on the size of $\frac{1} {|s||s_n|}$. The trick is to look far enough out into the sequence $\{s_n\}$ that the terms are closer to $s$ than they are to 0.

Consider the particular value $\varepsilon_0 = \frac{|s|}{2}$.  Because $\{s_n\} \to s$, there exists $N_1$ such that 
\[n \geq N_1 \implies |s_n - s|<\frac{|s|}{2}, \quad \text{and therefore} \quad |s_n| > \frac{|s|}{2}.\]
Next, we choose $N_2$ such that 
\[n \geq N_2 \implies |s_n -s| < \frac{\varepsilon |s|^2}{2}.\]

Let $N=\max\{N_1, N_2\}$, then $n \geq N$ implies that
\[\left\lvert \frac{1}{s_n} - \frac{1}{s} \right\rvert = |s-s_n|\frac{1} {|s||s_n|} < \frac{\varepsilon |s|^2}{2} \frac{1}{|s|\frac{|s|}{2}} = \varepsilon.\]



\end{proof}

\begin{theorem} Extending our previous results to sequences in $\mathbb{R}^k$

(a) Suppose $x_n \in \mathbb{R}^k$ $(n=1,2,3,\ldots)$ and
\[x_n = (\alpha_{1,n}, \ldots, \alpha_{k,n})\]
Then $\{x_n\}$ converges to $x = (\alpha_1, \ldots \alpha_k)$ iff 
\[\lim_{n\to\infty}\alpha_{j,n} = \alpha_j \quad (1\leq j \leq k)\]

(b) Suppose $\{x_n\}, \{y_n\}$ are sequences in $\mathbb{R}^k$, $\{\beta_n\}$ is a sequence of real numbers, and $x_n \to x$, $y_n \to y$, $\beta_n \to \beta$. Then
\[\lim_{n\to\infty} (x_n + y_n ) = x+y, \quad \lim_{n\to\infty} (x_n \cdot y_n ) = x\cdot y,  \quad \lim_{n\to\infty} \beta_n x_n  = \beta x\]
\end{theorem}

\begin{proof}

\end{proof}

The above theorem verifies that the relationship between algebraic combinations of sequences and the limiting process is as trouble-free as we could hope for.  Limits can be computed from the individual component sequences provided that each component limit exists.  The limiting process is also well-behaved with respect to the order operation.

\begin{theorem}[Order Limit Theorem]  Assume $(a_{n}) \to a$ and $(b_{n})\to b$
\begin{enumerate}[(i)]
	\item If $a_{n} \geq 0$ for all $n \in \mathbb{N}$, then $a\geq0$.
	\item If $a_{n}\leq b_{n}$ for all $n \in \mathbb{N}$, then $a\leq b$.
	\item If there exists $c \in \mathbb{R}$ for which $c \leq b_{n}$ for all $n \in \mathbb{N}$, then $c \leq b$.  Similarly, if $a_{n} \leq c$ for all $n \in \mathbb{N}$, then $a\leq c$.
\end{enumerate}
\end{theorem}

\begin{proof}
\begin{enumerate}[(i)]
	\item We will begin this proof by contradiction; that is, assume that $a<0$.  The idea is to produce a term in the sequence $(a_{n})$ that is also less than zero.  To do this, we consider a particular value $\varepsilon_{0}=|a|$. The definition of convergence guarantees that we can find an $N$ such that $|a_{n} - a|<|a|$ for all $n \geq N$.  In particular this would mean that $|a_{N} - a| < |a|$ which implies that $a_{N}<0$.  This contradicts our hypothesis that $a_{N} \geq 0$.  We therefore conclude that $a \geq 0$.
	
	\begin{marginfigure}
  \includegraphics{OrderLimitTH.jpg}
  \caption{If $a_{n} \geq 0$ for all $n \in \mathbb{N}$, then $a\geq0$.}
  %\label{fig:textfig1}
\end{marginfigure}
	\item  The Algebraic Limit Theorem ensures that the sequence $(b_{n} - a_{n})$  converges to $b-a$.  Because $b_{n} - a_{n} \geq 0$, we can apply the results from part (i) to get that $b-a \geq 0$, that is, $a \leq b$.
	\item  Take $a_{n} = c$ (or $b_{n} = c$) for all $n \in \mathbb{N}$ and apply our results from part (ii).

\end{enumerate}
\end{proof}


Loosely speaking, limits and their properties do not depend at all on what happens at the beginning of the sequence but are strictly determined by what happens as $n$ gets larger.  Changing the values of the first ten (or ten thousand) terms of the sequence has no effect on its eventual limit.


\medskip


\subsection{Exercises}
\noindent
\emph{Exercise 2.3.1.} Show that the constant sequence $(a,a,a,a,...)$ converges to $a$.
\begin{proof}
Let $\varepsilon > 0$.  Since the sequence is constantly $a$, we know that for all $n$:
	$$ |a_n - a| = |a-a| = 0 < \varepsilon $$
	Therefore we can choose $N$ completely arbitrarily since any choice will ensure that $|a_n -a| = 0 < \varepsilon$
\end{proof}

\smallskip

\noindent
\emph{Exercise 2.3.2.}
Let $x_n \geq 0$ for all $n \in \mathbb{N}$. 
\begin{enumerate}[(a)]
\item If $(x_n) \to 0$, show that $(\sqrt{x_n}) \to 0$. 
\item If $(x_n) \to x$, show that $(\sqrt{x_n}) \to \sqrt{x}$.
\end{enumerate}

\begin{proof}
\begin{enumerate}[(a)]
\item Let $\varepsilon> 0$.  (We want to find an $N$ such that $n \geq N$ implies $|\sqrt{x_n}| < \varepsilon$.)

Because $(x_n) \to 0$, there exists $N$ such that $n \geq N$ implies $|x_n|= x_n < \varepsilon^2$.  Then for the same $N$ and all $n \geq N$, $|\sqrt{x_n}| = \sqrt{x_n} < \sqrt{\varepsilon^2} = \varepsilon$.

As a result, for all $\varepsilon> 0$, there exists $N$ such that $n \geq N$ implies that $|\sqrt{x_n}- 0| < \varepsilon$ and $(\sqrt{x_n}) \to 0$.



\item The case $x = 0$ is part (a), so assume that $x > 0$.

Let $\varepsilon > 0$.  We wish to prove that there exists $N$ such that $n \geq N$ implies that $|\sqrt{x_n} - \sqrt{x}| < \varepsilon$.  

$\begin{array} {lcl} 
|\sqrt{x_n} -  \sqrt{x}|& = &|\sqrt{x_n} -  \sqrt{x}| \frac{\sqrt{x_n} +  \sqrt{x}}{\sqrt{x_n} +  \sqrt{x}}\\ 
	& = & \frac{|x_n -  x|}{\sqrt{x_n} +  \sqrt{x}}\\
&\leq	& \frac{|x_n -  x|}{\sqrt{x}}
    \end{array}$
    
    Because $(x_n ) \to x$ and $x>0$, we can choose $N$ such that $|x_n -x | <\varepsilon \sqrt{x}$ whenever $n \geq N$.  This implies that $\forall n \geq N$,
    \[		|\sqrt{x_n} - \sqrt{x}| \leq \frac{|x_n -  x|}{\sqrt{x}} < \frac{\varepsilon \sqrt{x}}{\sqrt{x}} = \varepsilon	\]

\end{enumerate} \end{proof}

\smallskip

\noindent
\emph{Exercise 2.3.3.} (Squeeze Theorem). Show that if $x_n \leq y_n \leq z_n$ for all $n \in \mathbb{N}$, and if $\lim x_n = \lim z_n = l$, then $\lim y_n = l$ as well.
\begin{proof}
Let $\varepsilon > 0$.  Since $(x_n) \to l$, there exists $N_{1}$ such that $n \geq N_{1} \implies x_n \in (l- \varepsilon, l + \varepsilon)$.

Since $(z_n) \to l$, there exists $N_{2}$ such that $n \geq N_{2} \implies z_n \in (l- \varepsilon, l + \varepsilon)$.

Since $x_n \leq y_n \leq z_n$ for all $n \in \mathbb{N}$, setting $N = \max \{N_1, N_2\}$ it follows that $y_n \in (l-\varepsilon, l + \varepsilon)$ for all $n \geq N$.

\end{proof}

\smallskip


\noindent
\emph{Exercise 2.3.4.}  Show that limits, if they exist, must be unique. In other words, assume $\lim a_n = l_1$ and $\lim a_n = l_2$, and prove that $l_1 = l_2$.
\begin{proof}
\begin{enumerate}[(i)]
	\item    Using the definition of convergence
	\noindent \marginnote{ \begin{proof} $\Rightarrow$ trivial.
		
		$\Leftarrow$ Suppose $l_1 \neq l_2$ and set $\varepsilon_{0} = |l_1 - l_2| > 0$.  Then the inequality $|l_1 - l_2| < \varepsilon$ fails for the particular choice $\varepsilon = \varepsilon_{0}$.  Contradiction.
		\end{proof}}\begin{lemma}  Two numbers $l_1$ and $l_2$ are equal iff $|l_1 - l_2 |< \varepsilon$ for all $\varepsilon > 0$.
		\end{lemma} 
	
	Let $\varepsilon > 0$.  Since $(a_n) \to l_1$ and $(a_n) \to l_2$, there exist $N_1, N_2$ such that $n \geq N_1$ and $n \geq N_2$ implies that $|a_n - l_1| < \varepsilon/2$ and $|a_n - l_2| < \varepsilon/2$  respectively.
	
	Therefore choosing $N = \max\{N_1, N_2\}, \quad n \geq N$ implies:


$\begin{array} {lcl} 
|l_1 - l_2|& = &|l_1 + a_n - a_n - l_2 | \\ 
	& \leq & |l_1 - a_n | +  |a_n - l_2 |\\
	& = & |a_n - l_1 | + |a_n - l_2 |\\
	& < & \frac{\varepsilon}{2} + \frac{\varepsilon}{2} \\
	&=& \varepsilon
    \end{array}$	 
    
    Since $\varepsilon$ was arbitrary, the lemma gives $l_1 = l_2$.

	\item  By the Algebraic Limit Theorem
	
	$\begin{array} {lclcl} 
\lim (a_n - a_n)& = &\lim (a_n) -\lim (a_n)  &=& l_{1} - l_{2}\\ 
			& = &\lim 0 			&=& 0 \\
    \end{array}	$
    
    Therefore $l_1 - l_2 = 0$ implying that $l_1 = l_2$.
	
\end{enumerate}
\end{proof}

\smallskip


\noindent
\emph{Exercise 2.3.6.} 
\begin{enumerate}[(a)]
\item Show that if $(b_n) \to b$, then the sequence of absolute values $(|b_n|)$ converges to $|b|$.


\item Is the converse of part (a) true? If we know that $|b_n| \to |b|$, can we deduce that $(b_n) \to b$?
\end{enumerate}

\begin{proof}
\begin{enumerate}[(a)]
\item We wish to show that $\forall \varepsilon > 0 \quad |b_n - b| < \varepsilon  \implies ||b_n| - |b|| < \varepsilon$.

$\begin{array} {rcrcl} 
|b_n|& = &|b_n - b + b|  &\leq & |b_n - b| + |b|\\ 
	&&	|b_n| - |b|	& \leq & |b_n - b|	\\
||b_n| - |b|| &\leq&    ||b_n - b|| &=& |b_n - b|\\
    \end{array}	$
    
    (The last line uses the same argument with the roles of $b_n$ and $b$ exchanged, which gives $|b| - |b_n| \leq |b_n - b|$ as well.)  Now let $\varepsilon > 0$.  Since $(b_n) \to b$, there exists $N$ such that $n \geq N \Rightarrow |b_n - b| < \varepsilon$.
    
    Therefore, for the same $N$ and all $n \geq N$,
    $$ ||b_n| - |b|| \leq     |b_n - b| < \varepsilon $$
    and $(|b_n|)$ converges to $|b|$.
    
\item No, the converse is not true.  Take $b_n = (-1)^n$: then $(|b_n|)$ converges to $|1|$, but $(b_n)$ does not converge to $1$ (or to anything else).
\end{enumerate}
\end{proof}

\smallskip


\noindent
\emph{Exercise 2.3.7.} 
\begin{enumerate}[(a)]
	\item Let $(a_n)$ be a bounded (not necessarily convergent) sequence, and assume $\lim b_n = 0$. Show that $\lim(a_n b_n ) = 0$. Why are we not allowed to use the Algebraic Limit Theorem to prove this?

\item Can we conclude anything about the convergence of $(a_n b_n)$ if we assume that $(b_n)$ converges to some nonzero limit $b$?
%\item Use $(a)$ to prove Theorem 2.3.3, part (iii), for the case when $ a = 0$.
\end{enumerate}

\begin{proof}

\begin{enumerate}[(a)]
	\item   Let $\varepsilon > 0$.  Because $(a_n)$ is bounded, there exists $M > 0$ such that $|a_n| \leq M$ for all $n \in \mathbb{N}$.
	
	Furthermore since $(b_n) \to 0$, there exists $N$ such that $n \geq N$ implies that $|b_n| < \frac{\varepsilon}{M}$.
	
	Then the sequence $(a_n b_n) $ is such that for $n \geq N$
	$$ |a_n b_n | \leq |a_n| |b_n| < M \frac{\varepsilon}{M} = \varepsilon $$
	
	
	and $a_n b_n  \to 0$
	
	
	We cannot use the Algebraic Limit Theorem to prove this since $(a_n)$ is not assumed to converge.
	
	
	

\item No, we cannot conclude anything about the sequence.  With $b_n = 1$ for every $n$ the product is just $(a_n)$ itself, which may converge or diverge. 
\end{enumerate}
\end{proof}

\noindent
\emph{Exercise 2.3.10.}
If $(a_n) \to 0$ and $|b_n - b| \leq a_n$ for all $n$, then $(b_n) \to b$.

\begin{proof}
	Since $0 \leq |b_n - b| \leq a_n$, every $a_n$ is nonnegative.  Therefore $|a_n| = a_n$.
	
	Let $\varepsilon > 0$.  Since $(a_n) \to 0$, there exists $N$ such that $n \geq N$ implies $|a_n | = a_n < \varepsilon$.
	
	Then the inequality above gives $|b_n - b| \leq a_n < \varepsilon$ whenever $n \geq N$.   Therefore $(b_n) \to b$.
\end{proof}


\noindent
\emph{Exercise 14.}

  If $\{s_{n}\}$ is a complex sequence, define its arithmetic means $\sigma_{n}$ by
 $$  \sigma_{n} = \frac{s_{0} + s_{1} + \cdots + s_{n}}{n+1} \qquad (n=0,1,2,\ldots) $$
 
 \begin{enumerate}[(a)]
 	\item  If $\lim s_{n} = s$, prove that $\lim \sigma_{n} = s$.
	
	\begin{proof}
			Let $\varepsilon >0$.
			
			If $\lim s_{n} = s$  then the sequence is bounded, that is, there exists an $M >0$ such that $$ |s_n - s| < M \quad \text{for all } n. $$
			
			Furthermore, there exists an $N_{1} \in \mathbb{N}$ such that $n \geq N_{1}$ implies that $$  |s_n - s|< \frac{\varepsilon}{2} $$
			
			Therefore for $n \geq N_{1}$ we have that 
			
			\begin{eqnarray*}
|\sigma_n - s|  &=& \big|\frac{s_0 + s_1 + s_2 + \cdots + s_n}{n+1}  - s\big|  \\
   	&=& \big|\frac{s_0 + s_1 + s_2 + \cdots + s_n - (n+1)s}{n+1}\big| \\
   &=& \big|\frac{(s_0-s) + (s_1 -s) + (s_2 -s)  + \cdots + (s_n - s)}{n+1}\big| \\
   &\leq & \big|\frac{(s_0-s) + (s_1 -s) +\cdots + (s_{N_{1} - 1} -s) }{n+1}\big|  + \big|\frac{(s_{N_{1}} -s) + \cdots + (s_n - s)}{n+1}\big| \\
   &<& \frac{N_{1} M}{n+1}  + \frac{\varepsilon}{2}\frac{(n - N_{1} + 1)}{n+1}\\
   &<& \frac{N_{1} M}{n+1}  + \frac{\varepsilon}{2}
\end{eqnarray*}

Note that in the last line we used the fact that for $n \geq N_1$, $\frac{(n - N_{1} + 1)}{n+1} < 1$. 


But $N_{1}$ and $M$ are fixed meaning they are simply constants.  Since the sequence $\{\frac{1}{n+1}\} \rightarrow 0$ we can make the quantity $\frac{N_{1} M}{n+1}$ to be arbitrarily small.

Choose $N_2$ such that $n \geq N_2$ implies that $$ \frac{N_{1} M}{n+1} < \frac{\varepsilon}{2} $$




Finally, let $N = \max\{N_1, N_2\}$.   Then $n \geq N$ implies that 


\begin{eqnarray*}
|\sigma_n - s|  &<& \frac{N_{1} M}{n+1}  + \frac{\varepsilon}{2}\\
		&<& \frac{\varepsilon}{2} + \frac{\varepsilon}{2} \\
		&=&  \varepsilon
\end{eqnarray*}

as desired.	
	\end{proof}
	\item  Construct a sequence $\{s_{n}\}$ which does not converge although $\lim \sigma_{n} = 0$.
	
	\noindent
	\emph{Solution:}   For any $a \neq 0$, the sequence $\{s_{n}\} = (a, -a, a, -a, \ldots )$ does not converge, but the sequence $\{\sigma_{n}\}$ converges to $0$.
 \end{enumerate}

 


\smallskip

\section{The Monotone Convergence Theorem }
We showed previously that convergent sequences are bounded.  While the converse is not necessarily true, if a sequence is bounded as well as monotonic, then it converges.


\begin{definition}Increasing sequence:  A sequence $\{a_{n}\}$ is \emph{increasing} if $a_{n} \leq a_{n+1}$ for all $n \in \mathbb{N}$ and \emph{decreasing} if $a_{n} \geq a_{n+1}$ for all $n \in \mathbb{N}$.  A sequence is \emph{monotone} if it is either increasing or decreasing. \end{definition}

\begin{theorem}  [The Monotone Convergence Theorem] If a sequence is monotone and bounded, then it converges.
\end{theorem}

\begin{proof}
	Let $(a_n)$ be monotone and bounded.
	WLOG assume $(a_n)$  is increasing and consider the set of points $\{a_n : n \in \mathbb{N}\}$.
	
	By assumption, the set is bounded.  By the axiom of completeness, the set has a supremum.  Let $s = \sup\{a_n : n \in \mathbb{N}\}$.
	
	We claim that $\lim (a_n) = s$.
	
	Let $\varepsilon > 0$.  Because $s$ is the least upper bound of $\{a_n : n \in \mathbb{N}\}$, $s- \varepsilon$ is not an upper bound, so there exists  a point in the sequence call it $a_N$ such that $s-\varepsilon < a_N$.
	
	Now the fact that $(a_n)$ is increasing implies that if $n \geq N$ then $a_N \leq a_n$.
	
	Hence,
	$$ s- \varepsilon < a_N \leq a_n \leq s < s+\varepsilon $$
	which implies that $|a_n  - s| < \varepsilon$ for all $n \geq N$.
\end{proof}

Note that the Monotone Convergence Theorem is extremely useful since it asserts the convergence of a sequence without explicit mention of the actual limit.

Also note that the Monotone Convergence Theorem follows from the Axiom of Completeness.

\subsection{Exercises}

	\emph{Exercise 2.4.2.}
		\begin{enumerate}[(a)]
			\item Prove that the sequence defined by $x_1 = 3$ and 
				\[		x_{n+1} = \frac{1}{4-x_n}\]
				converges.
			\item Now that we know $\lim x_n$ exists, explain why $\lim x_{n+1}$ must also exist and equal the same value.
			\item Take the limit of each side of the recursive equation in part (a) of this exercise to explicitly compute  $\lim x_n$.
		\end{enumerate}
		
		\begin{proof}
			\begin{enumerate}[(a)]
				\item  
				\begin{enumerate}[(i)]
					\item The first few terms are $3, 1, \frac{1}{3}, \frac{3}{11}, \ldots$.  Then we can guess that the sequence is decreasing ($x_{n+1}< x_{n} \quad \forall n$)  and we can prove it by induction.
				
				For $n=1$,  $x_{2}< x_{1}$ since $1 <3$.
				
				For $n=k$, assume that $x_{k+1} < x_{k}$.
				
				Then when $n=k+1$, we want to be able to show that $x_{k+2} =\frac{1}{4-x_{k+1}} < x_{k+1} = \frac{1}{4-x_{k}}  $.   Using the inductive hypothesis, we have that
				
				
				\begin{eqnarray*}
x_{k+1}  &< & x_{k} \\
   	-x_{k+1}  &>& -x_{k} \\
 4-x_{k+1}  &>&4 -x_{k}\\
    \frac{1}{4-x_{k+1}}  &<&\frac{1}{4 -x_{k}}\\
   x_{k+2}  &< & x_{k+1} \\
\end{eqnarray*}

as desired.
				\item We can see that $(x_n)$ is bounded above by $3$ but we would like to prove that it is bounded below by $0$.  Again we can prove this by induction, showing that $x_n > 0 \quad \forall n$.
				
				Clearly for $n=1$, $x_1 = 3 >0$.
				Then assume for $n=k$, that $x_k >0$.
				
				To show this for $n = k+1$, note that we have that $3>x_k>0$.  This implies that $x_{k+1} = \frac{1}{4-x_k}$  is positive.  Therefore $x_{k+1}>0$ as desired.
				
				
				Since $(x_n)$ is decreasing and bounded below by $0$, we have by the MCT that the sequence is convergent.
				
								\end{enumerate}
				
				\item Since $ (x_n)$ is bounded and decreasing, $(x_{n+1})$ must also be bounded and decreasing, then by the MCT it must also be convergent.
				
				Furthermore since $(x_n) \to L$, it must be that for every $\varepsilon >0$ there exists $N \in \mathbb{N}$ such that $n \geq N$ implies that $|x_n - L| <\varepsilon$.  However, notice that for $(x_{n+1})$ this same $N$ works since $n+1> n \geq N$ and therefore $|x_{n+1} - L| <\varepsilon$ and both sequences converge to the same limit.
				
							\item By part (b), we have that  $\lim x_n = \lim x_{n+1} = L$. Furthermore, by the Algebraic Limit Theorem we can deduce that $L = \lim x_{n+1} = \lim \frac{1}{4-x_n} = \frac{1}{4-L}$ and we have that $L$ satisfies the equation $L^2 - 4L +1 =0$.  Solving for $L$ we have that $L = 2 \pm \sqrt{3}$, however since $3$ is an upper bound for the sequence, we must have that $L= 2 - \sqrt{3}$.
			\end{enumerate}
		
		\end{proof}
		
\noindent		
\emph{Exercise 2.4.3.} Using the same methods, prove that the following sequence is convergent and find its limit.
	$$y_{n+1} = 4 - \frac{1}{y_n}	\qquad y_1 = 1$$
	
	\begin{proof}
		\begin{enumerate}[(i)]
			\item Monotone.  Looking at the first few terms of the sequence again, we can see that it is increasing: $y_n \leq y_{n+1}$ for all $n$.  To prove this we use induction (together with the fact, shown in part (ii), that every term is positive).
				
				For $n=1$, $y_1 \leq y_2$ since $1 < 3$.
				
				Assume that for $n=k$, it is true that $y_{k} \leq y_{k+1}$.
				
				Then for $n = k+1$, 
				\begin{eqnarray*}
y_{k}  &\leq & y_{k+1} \\
   	\frac{1}{y_{k}}  &\geq& \frac{1}{y_{k+1}} \\
    -\frac{1}{y_{k}}  &\leq& -\frac{1}{y_{k+1}}\\
      4- \frac{1}{y_{k}}  &\leq&4- \frac{1}{y_{k+1}}\\
   y_{k+1}  &\leq& y_{k+2} \\
\end{eqnarray*}

as desired.

\item Bounded.  We show by induction that $1 \leq y_n < 4$ for every $n$.  For $n=1$ this is trivial.  Assume then that $1 \leq y_k < 4$.  Then $0 < \frac{1}{y_k}\leq 1$, and so $3 \leq y_{k+1} = 4 - \frac{1}{y_k} < 4$ as well.

\item Limit.  By the MCT the sequence converges; call its limit $y$, so that $\lim y_n = \lim y_{n+1} = y$.  Then $y = \lim y_{n+1} =  \lim (4 - \frac{1}{y_n}) = 4 - \frac{1}{y}$.  By this relation we have that $y^2 - 4y + 1 =0$ with roots $y = 2 \pm \sqrt{3}$.  However, since the sequence is bounded below by 1, we have that $y = 2 + \sqrt{3}$.

				
		\end{enumerate}
	\end{proof}
	
\emph{Exercise 2.4.4.}  Prove that the sequence $x_n = \sqrt{2}, \sqrt{2\sqrt{2}}, \sqrt{2\sqrt{2\sqrt{2}}}, \ldots$
	converges and find the limit.
	
	\begin{proof}
	First, define the sequence recursively:
	$$  x_{n+1} = \sqrt{2 x_n} \qquad x_1 = \sqrt{2} $$
		\begin{enumerate}[(i)]
		
			\item Monotone Increasing  ($x_{n} < x_{n+1}\quad \forall n$).  For $n=1, \quad \sqrt{2} < \sqrt{2\sqrt{2}}$. Then assume that for $n=k$, $x_k < x_{k+1}$.
			
			Using the Inductive hypothesis
			\begin{eqnarray*}
x_{k}  & <& x_{k+1} \\
   	2x_{k}  &<& 2x_{k+1} \\
    \sqrt{2x_{k}}  &<&\sqrt{ 2 x_{k+1}}\\
    x_{k+1}  & <& x_{k+2}
\end{eqnarray*} as desired.

\item Bounded.  We wish to prove that the sequence is bounded above by 2.
For $x_1$ this is trivial.  Assume that $x_k < 2$.
Then 
\marginnote{Notice that here we take the limit under a square root which is not covered in the Algebraic Limit Theorem.  We therefore use Exercise 2.3.2 to support our argument.}
\begin{eqnarray*}
x_{k}  & <& 2 \\
   	2x_{k}  &<& 4 \\
    \sqrt{2x_{k}}  &<&\sqrt{ 4}\\
 & <& 2 \\
 x_{k+1}  & <& 2
\end{eqnarray*} as desired.   


\item Limit.  Note that $x = \lim x_n = \lim x_{n+1}$.  Therefore we have that $x = \lim x_{n+1} = \sqrt{2x}$ and we have that $x^2 - 2x = 0$, so $x = 0$ or $x = 2$.  Since the sequence increases from $x_1 = \sqrt{2} > 0$, the limit cannot be $0$.  Therefore $x = 2$.
		\end{enumerate}
	\end{proof}
	
	\emph{Exercise 2.4.6.} Let $(a_n)$ be a bounded sequence.
	
	\begin{enumerate}[(i)]
	\item Define the sequence $y_n = \sup \{a_k : k \geq n\}$.  Prove that $\lim y_n$ exists.  This limit is called the \emph{limit superior} of $(a_n)$ and is written $\limsup a_n$.
	\item Define the sequence $z_n = \inf \{a_k : k \geq n\}$ and prove that it converges.  Its limit is the \emph{limit inferior}, written $\liminf a_n$.
	\item Show that $\limsup a_n = \liminf a_n$ if and only if $\lim a_n$ exists, in which case all three values are equal.
	\end{enumerate} % COMPLETE


\section{Subsequences and Bolzano-Weierstrass Theorem}

\begin{definition}Subsequences:  Given a sequence $\{p_n\}$ consider a sequence $\{n_k\}$ of positive integers such that $n_1 < n_2 < n_3 < \cdots$.  Then the sequence $\{p_{n_k}\}$ is called a \emph{subsequence} of $\{p_n\}$.  If $\{p_{n_k}\}$ converges, its limit is called a \emph{subsequential limit} of $\{p_{n}\}$.
\end{definition}
\medskip
Note:  It is clear that $\{p_{n}\}$ converges to $p$ if and only if every subsequence of $\{p_{n}\}$ converges to $p$. 

Notice also that the order of the terms in the subsequence is the same as in the original sequence and that repetitions are not allowed.

\begin{theorem} Subsequences of a convergent sequence converge to the same limit as the original sequence.\end{theorem}

\begin{proof} 
Assume the sequence $(x_n) \to L$ and let $(x_{n_j})$ be a subsequence of this sequence.  
Let $\varepsilon > 0$.  Then there exists an $N$ such that $n \geq N$ implies that $|x_n - L| < \varepsilon$.  But since $n_j \geq j$ this same $N$ works for the subsequence as well.  To be more precise, $j\geq N$ implies that $n_j \geq N$ and that $|x_{n_j} - L| < \varepsilon$ as desired.


\end{proof}

This theorem can sometimes provide a very useful proof for divergence of a sequence.  If we can find two subsequences that converge to two different limits, then we can rigorously conclude that the original sequence diverges.

Furthermore, this theorem can be applied to infinite series to show that infinite sums are associative which we will see later.

\begin{example} Let $0 < b< 1$.  Because $b >b^2 > b^3 > b^4 > \ldots > 0$ the sequence $(b^n)$ is decreasing and bounded below.  By the MCT the sequence converges to some limit $l$ satisfying $b > l \geq 0$.   Notice however that $(b^{2n})$ is a subsequence, so by the above theorem $(b^{2n}) \to l$. But $(b^{2n}) = (b^n)( b^n)$ so by the Algebraic limit theorem, $(b^{2n}) \to l \cdot l = l^2$.  Because limits are unique, $l^2 = l$, and since $l \leq b < 1$ rules out $l = 1$, we conclude that $l=0$.
\end{example}

\begin{example}
Consider the sequence 
$$\big( 1, - \frac{1}{2}, \frac{1}{3}, -\frac{1}{4}, \frac{1}{5}, -\frac{1}{5}, \frac{1}{5}, -\frac{1}{5}, \frac{1}{5}, -\frac{1}{5},\ldots	\big)$$

We can show that the sequence is divergent since we can provide two subsequences, namely:
$$ \big( \frac{1}{5}, \frac{1}{5}, \frac{1}{5}, \frac{1}{5},  \ldots \big) \quad \text{and} \quad \big( -\frac{1}{5},-\frac{1}{5},-\frac{1}{5},-\frac{1}{5}, \ldots \big) $$

that converge to two different limits.
\end{example}

\begin{theorem}\newthought{Bolzano-Weierstrass Theorem:} Every bounded sequence contains a convergent subsequence.\end{theorem}

\begin{proof}
Let $(a_n)$ be a bounded sequence such that $|a_n| \leq M$ for all $n$.  Bisect the closed interval $[-M, M]$ into $[-M, 0]$ and $[0,M]$ (the midpoint belongs to both halves).  Now it must be that at least one of these closed intervals contains infinitely many terms of the sequence $(a_n)$.  Select the half for which this is the case and label it $I_1$.  Then let $a_{n_1}$ be some point in the sequence satisfying $a_{n_1} \in I_1$.

Now bisect $I_1$ and again select the half with infinitely many terms of $(a_n)$ and call it $I_2$.  Then again select a point $a_{n_2} \in I_2$ such that $n_2> n_1$.

Continue in this way, constructing $I_{k}$ by halving $I_{k-1}$ and choosing the half with infinitely many terms of the sequence $(a_n)$ and then selecting $a_{n_k} \in I_{k}$ such that
\[	n_k  > n_{k-1} > n_{k-2} > \cdots > n_{2} > n_{1}	\]


We want to show that $(a_{n_k})$ converges but we first need to find a candidate for the limit.

Note that 
\[	I_1 \supseteq I_2 \supseteq I_3 \supseteq I_4 \supseteq	\cdots \]
and that each $I_i$ is closed.

By the NIP, there exists $x \in \mathbb{R}$ such that $x \in I_k$ for every $k$.


Let $\varepsilon >0$.  By construction, the length of $I_k$ is $M (\frac{1}{2})^{k-1}$, which tends to $0$ as $k \to \infty$.

Therefore there exists $N$ such that $k \geq N$ implies $M (\frac{1}{2})^{k-1} < \varepsilon$.  But $x$ and $a_{n_k}$ both lie in $I_k$, so the distance between them is at most the length of $I_k$.  As a result, $|x-a_{n_k}|< \varepsilon$ for all $k \geq N$, and $(a_{n_k}) \to x$.

\end{proof}

\section{The Cauchy Criterion}  
\newthought{Cauchy Sequence:}  A sequence $(a_n)$ is called a \emph{Cauchy sequence} if, for every $\varepsilon >0$ there exists an $N \in \mathbb{N}$ such that whenever $m,n \geq N$ it follows that $|a_n - a_m| < \varepsilon$

That is, a sequence is Cauchy if, for every $\varepsilon$, there is a point in the sequence after which all the terms are closer to each other than the given $\varepsilon$.  To spoil the surprise, we will argue that Cauchy sequences are convergent and that convergent sequences are Cauchy.  This, like the Monotone Convergence Theorem, allows us to prove that a sequence converges without having explicit knowledge of what the limit might be. 


\begin{lemma} Cauchy Sequences are bounded.\end{lemma}

\begin{proof}
	Given $\varepsilon = 1$, there exists an $N$ such that $|x_m - x_n| <1$ for all $m,n \geq N$.  
	As a result \marginnote{Note that $|x_n| = |x_n + x_N - x_N| \leq |x_n -x_N| + |x_N| < 1 + |x_N|$}  $|x_n| < |x_N| + 1$ for all $n \geq N$. Note that there were finitely many terms before $x_N$ in the sequence. Therefore let
	
	
	\[	M	= \max \{|x_1|,|x_2|, \ldots, |x_{N-1}|, |x_N| + 1 \}\]
	And this value is a bound for the sequence.
\end{proof}


\begin{theorem}\newthought{Cauchy Criterion:} A sequence converges iff it is a Cauchy sequence. \end{theorem}

\begin{proof}
	$\Longrightarrow$.  This we can prove by applying the triangle inequality. Assume that a sequence $(x_n)$ converges to a limit $x$.   Then $\forall \varepsilon >0$, there exists an $N$ such that $n,m \geq N$ implies that 
	\[ |x_n -x | < \varepsilon/2	\qquad 	|x_m - x| < \varepsilon/2	\]
Applying the triangle inequality, we can state that 

$$
\begin{array}{rcl}
|x_n - x_m| &=& |x_n - x  + x - x_m| ,\\
&\leq & |x_n - x| + |x - x_m|,\\
& < & \varepsilon\\
\end{array}
$$

as desired.   

	$\Longleftarrow$.  So the other direction is a lot harder to prove and we have to invoke Bolzano-Weierstrass to get our result.
	
	Assume we have a Cauchy sequence $(x_n)$.  By the lemma above we know that it is bounded and so by BW we know that it contains a convergent subsequence.   Let us call this subsequence $(x_{n_k})$ and call its limit $x$.  
	
	The idea here is that because the sequence is Cauchy, we know that elements in the ``tail'' end of the sequence are getting closer together.  But because the terms in the subsequence are getting closer to the limit $x$, we can force each value to be less than $\varepsilon/2$ and apply the triangle inequality.
	
	So let $\varepsilon >0$.
	
		Because $(x_n)$ is Cauchy there is an $N$ such that $m,n \geq N$ implies that \[|x_n - x_m|<\varepsilon/2.\]  We also have that the subsequence converges so we can choose an index, let us say $n_K$, such that $n_K \geq N$ and 
		\[	|x_{n_K} - x|<\varepsilon/2 \]
		
		
		Now take any $n \geq N$.  Since $n_K \geq N$ as well, the Cauchy condition applies to the pair $n$, $n_K$, and so 
		
				$$
				\begin{array}{rcl}
|x_n - x| &=& |x_n - x_{n_K}  + x_{n_K} - x| ,\\
&\leq & |x_n - x_{n_K}| + |x_{n_K} - x|,\\
& < & \varepsilon
\end{array}
$$
	
\end{proof}



One important thing to note is that many of our significant results up to this point rely heavily on the Axiom of Completeness.

The AoC was what allowed us to prove the NIP which allowed us to prove Bolzano-Weierstrass which allowed us to prove the Cauchy Criterion.  Furthermore the AoC also gave us the Monotone Convergence Theorem.

$$
\begin{array}{rcl}
AoC &\implies & NIP \implies BW \implies CC ,\\
&\implies & MCT,\\
\end{array}
$$

Also note that we could have begun this chain with \emph{any} of these statements and arrived at the others as a result.\marginnote{A subtlety: starting from the Nested Interval Property or the Cauchy Criterion, one must also assume the Archimedean Property in order to recover the Axiom of Completeness.}  This is pretty remarkable.  In other words all of these statements are equivalent and that assuming any one of them to be true, means that we can derive the others.   This is because all of these statements essentially say the same thing : that the field of real numbers is \emph{complete}.   Remember how we could not show the AoC in the field of rational numbers.  This is because the rationals did not contain a certain property, independent of the field and order axioms, that forced this result to be true.  In the next section we try to define more explicitly what this ``completeness'' property is and how it relates to the topology of the field in question.


\section{Infinite Series}


\newthought{Infinite Series:}  Let $(b_n)$ be a sequence.  An infinite series is an expression of the form
		\[\sum_{n=1}^{\infty} b_n = b_1 + b_2 + b_3 + \cdots \]
		
		The sequence of \emph{Partial sums} of the series $(s_m)$ is defined as 
			\[s_m = b_1 + b_2 + b_3 + \cdots + b_m \]
			
		We say that the series $\sum_{n=1}^{\infty} b_n$ \textsc{converges} if its sequence of partial sums converges to some limit B and we denote this limit as $\sum_{n=1}^{\infty} b_n  = B$.
		Specifically,
		\[\sum_{n=1}^{\infty} b_n = B \qquad \text{means that} \qquad \lim s_n = B\]
		
		
		\begin{dangerousbend}
			It is important to keep a distinction between 
				\begin{itemize}
					\item The sequence of \emph{terms}: $(b_1, b_2, b_3, \ldots)$ and
					\item The sequence of \emph{partial sums}:  $(s_1, s_2, s_3, \ldots)$ where $s_n = b_1 +  \cdots + b_n$
				\end{itemize}
		\end{dangerousbend}
		
		Nonetheless it is this close relationship between series and sequences that allows us to apply much of our previous knowledge to them.
		
		
		\begin{theorem}\newthought{Theorem: Algebraic Limit Theorem for Series}
		If $\sum_{n=1}^{\infty} a_n = A$ and $\sum_{n=1}^{\infty} b_n = B$ then,
			\begin{enumerate}[(1)]
				\item $\sum_{n=1}^{\infty} c a_n = cA$
				\item $\sum_{n=1}^{\infty} (a_n + b_n) = A + B$
			\end{enumerate}
		\end{theorem}
		
		\begin{proof}
		\end{proof}
		
		\begin{theorem} \newthought{Cauchy Criterion for Series}
		The series $\sum_{n=1}^{\infty} a_n$ converges 
		
		\[\text{if and only if}\]
		
		
		 for every $\varepsilon > 0$ there exists an $N \in \mathbb{N}$ such that $n>m\geq N$ implies that $|a_{m+1} + a_{m+2} + \cdots + a_n| < \varepsilon$
		\end{theorem}
		
		
	\begin{proof}  \end{proof}
	
	This can help us prove some basic facts about series.
	
	
	\begin{theorem} If the series $\sum_{n=1}^{\infty} a_n$ converges then $(a_k) \to 0$.  \marginnote{The converse statement is \emph{not} true! Take for example the harmonic series.}
	\end{theorem}
	
	\begin{proof}  \end{proof}
	
	\begin{theorem}\newthought{The Comparison Test}  Assume that $(a_n)$ and $(b_n)$ are sequences where $0 \leq a_n \leq b_n$  for all $n \in \mathbb{N}$
		\begin{enumerate}
			\item If $\sum_{n=1}^{\infty} b_n$ converges, then $\sum_{n=1}^{\infty} a_n$ converges.
			\item If $\sum_{n=1}^{\infty} a_n$ diverges, then $\sum_{n=1}^{\infty} b_n$ diverges.
		\end{enumerate}
	\end{theorem}
	
	
	
	\begin{proof}  \end{proof}
	
	Remember also that statements about the convergence of a sequence (or series) only care about the tail ends of the sequence.  In other words, we can change around a few things about some finite number of initial terms without affecting the infinite behavior of the sequence (or series).
	
	Because the Comparison Test is about deducing the behavior of a series in comparison to another, it might be a good idea to have an arsenal of basic series to which we can compare another. 
	
	\newthought{Absolute Convergence:}  If $\sum_{n=1}^{\infty} |a_n|$ converges then we say that $\sum_{n=1}^{\infty} a_n$ converges absolutely.  On the other hand if $\sum_{n=1}^{\infty} a_n$ converges and $\sum_{n=1}^{\infty} |a_n|$ diverges we say that the series \emph{converges conditionally}.
	
	
	\begin{theorem}\newthought{Absolute Convergence Test}
	If the series $\sum_{n=1}^{\infty} |a_n|$ converges, then $\sum_{n=1}^{\infty} a_n$ converges as well.
	\end{theorem}  \marginnote{The converse of this theorem is also false!}
	\begin{proof}  \end{proof}
	
	
	\begin{theorem}\newthought{Alternating Series Test}  Let $(a_n)$ be a sequence satisfying,
	
	\begin{enumerate}
		\item The sequence is monotone decreasing:  $a_1 \geq a_2 \geq a_3 \geq \cdots a_n \geq a_{n+1} \geq \cdots $ and
		\item $(a_n) \to 0$
	\end{enumerate}
	
	
	Then the alternating series $\sum_{n=1}^{\infty} (-1)^{n+1} a_n$ converges.
	\end{theorem}
	\begin{proof}  \end{proof}
	
	\begin{theorem}\newthought{Ratio Test}
	Given a series $\sum_{n=1}^{\infty} a_n$ with $a_n \neq 0$, the ratio test states that if $(a_n)$ satisfies 
	\[		\lim \left|\frac{a_{n+1}}{a_n}	\right| = r < 1\]
	then the series converges absolutely (and therefore converges).
	\end{theorem} 
	\begin{proof}  \end{proof}
	
	
	\subsection{Rearrangements}
		Imagine that we were to permute the terms in a sum into some other order such that all of the original terms eventually appeared in the new order and such that no terms were repeated.  Clearly there are an infinite number of possible rearrangements of a series.  
		
		More formally,
		
		\newthought{Rearrangement}  Let  $\sum_{n=1}^{\infty} a_n$ be a series.  Another series  $\sum_{n=1}^{\infty} b_n$ is called a rearrangement of  $\sum_{n=1}^{\infty} a_n$ if there exists a bijection $f: \mathbb{N} \to \mathbb{N}$ such that $b_{f(k)} = a_{k}$ for all $k$.
		
		It should be noted however, that rearrangements of series do \emph{not} always converge to the same limit as the original series.  This is only true for absolutely convergent series as we will now prove.
		
		
		\begin{theorem}
	If $\sum_{n=1}^{\infty} a_n$ converges absolutely, then any rearrangement of the series converges to the same limit.
	\end{theorem} 
	\begin{proof}  \end{proof}
	
	This theorem should make clear that absolute convergence is an extremely desirable property when manipulating series.  However it should be noted that conditional convergence has this amazing and not at all obvious property which we attribute to Riemann.
	
	
	\begin{theorem}\newthought{Riemann}
	If $\sum_{n=1}^{\infty} a_n$ converges conditionally, then for any real number $r$ we can find a rearrangement of the series that converges to $r$.
	\end{theorem} 
	\marginnote{This is Theorem 3.54 in Rudin's \emph{Principles of Mathematical Analysis}.}
		
		


\section{Problems}

\noindent The problems in this section are original to this text.

\begin{problem}[Three proofs on trial]
\hint{Read each proof against the template in Section~\ref{sec:sequence-examples}.  Who is supposed to speak first, $\varepsilon$ or $N$?  And for which $n$ must the final inequality hold?}
Each of the following ``proofs'' has something wrong with it.  In each case say whether the \emph{claim} is true, pinpoint the faulty step, and, where the claim is true, write a correct proof.
\begin{enumerate}[(a)]
	\item \emph{Claim.} $\lim \frac{1}{n} = 0$.

	\emph{``Proof.''}  Let $N \in \mathbb{N}$ be arbitrary and set $\varepsilon = 2/N$.  If $n \geq N$ then $\left|\frac{1}{n} - 0\right| = \frac{1}{n} \leq \frac{1}{N} < \frac{2}{N} = \varepsilon$.  Hence $\left(\frac{1}{n}\right) \to 0$.
	\item \emph{Claim.} the sequence $a_n = (-1)^n$ converges to $1$.

	\emph{``Proof.''}  Let $\varepsilon > 0$ and choose $N = 2$.  For every even $n \geq N$ we have $|a_n - 1| = |1 - 1| = 0 < \varepsilon$.
	\item \emph{Claim.} $\lim \frac{2n+1}{n} = 2$.

	\emph{``Proof.''}  Let $\varepsilon > 0$ and choose $N = 1/\varepsilon$.  If $n \geq N$ then $\left|\frac{2n+1}{n} - 2\right| = \frac{1}{n} \leq \frac{1}{N} = \varepsilon$.
\end{enumerate}
To convince yourself that (a) is a broken argument and not merely an unusual one, show that exactly the same reasoning ``proves'' that the constant sequence $(1, 1, 1, \ldots)$ converges to $0$.
\end{problem}

\begin{problem}[Sticking and visiting]
\hint{For the second half of (c), look back at Chapter One and ask how many rational numbers there are.  For (d), the Bolzano--Weierstrass Theorem will find you a second number to visit.}
Here are two invented definitions, built from the same ingredients as convergence but with the quantifiers rearranged.
\begin{itemize}
	\item A sequence $(a_n)$ \emph{sticks to} $a$ if there exists an $N \in \mathbb{N}$ such that for every $\varepsilon > 0$ and every $n \geq N$ we have $|a_n - a| < \varepsilon$.
	\item A sequence $(a_n)$ \emph{visits} $a$ if for every $\varepsilon > 0$ and every $N \in \mathbb{N}$ there exists an $n \geq N$ with $|a_n - a| < \varepsilon$.
\end{itemize}
\begin{enumerate}[(a)]
	\item Describe each definition in a plain sentence that mentions neither $\varepsilon$ nor $N$.
	\item Show that if $(a_n)$ sticks to $a$ then it converges to $a$, and that if it converges to $a$ then it visits $a$.  Give examples to show that neither implication can be reversed.
	\item Give an example of a sequence that visits exactly two numbers.  Then give an example of a single sequence that visits \emph{every} real number.
	\item Suppose $(a_n)$ is bounded, visits $a$, and visits no other number.  Prove that $(a_n) \to a$.  Show by example that the conclusion can fail if the word ``bounded'' is dropped.
\end{enumerate}
\end{problem}

\begin{problem}[Climbing toward a hole]
\hint{For (c), write $x_{n+1} - x_n$ as a single fraction and look at the sign of the numerator.}
Define a sequence by $x_1 = 1$ and
\[ x_{n+1} = \frac{2x_n + 3}{x_n + 2}. \]
\begin{enumerate}[(a)]
	\item Compute $x_2$, $x_3$ and $x_4$, both as fractions and as decimals.  Explain why every $x_n$ is a positive rational number.
	\item Verify the identity
	\[ 3 - x_{n+1}^2 = \frac{3 - x_n^2}{(x_n + 2)^2} \]
	and use it to prove that $x_n^2 < 3$ for every $n$.
	\item Show that $(x_n)$ is increasing, and conclude that it converges.  Find the square of its limit.
	\item Every term of this sequence lies in $\mathbb{Q}$, and within $\mathbb{Q}$ the sequence is increasing and bounded.  Which theorem of this chapter have you just shown to be false in $\mathbb{Q}$?
	\item Use the identity in (b) to show that $3 - x_n^2 \leq 2/9^{\,n-1}$.  How many terms are needed before $x_n^2$ is guaranteed to be within one millionth of $3$?
\end{enumerate}
Notice what part (c) accomplishes.  Without ever writing the symbol $\sqrt{3}$, you have proved that there is a positive real number whose square is $3$.
\end{problem}

\begin{problem}[Bolzano--Weierstrass by decimal digits]
\hint{Part (a) is the pigeonhole principle for infinitely many pigeons.  If infinitely many terms are shared among ten bins, the bins cannot all be finite.}
The proof of the Bolzano--Weierstrass Theorem in the text cuts an interval in half again and again.  This problem redoes it by cutting into tenths, which lets us watch the limit being written down one digit at a time.  Let $(a_n)$ be any sequence with $0 \leq a_n < 1$ for all $n$.
\begin{enumerate}[(a)]
	\item Sort the terms into ten bins $[0, 0.1)$, $[0.1, 0.2)$, \ldots, $[0.9, 1)$ according to their first decimal digit.  Explain why at least one bin receives $a_n$ for infinitely many values of $n$.  Call the digit of one such bin $d_1$.
	\item Split that bin into ten smaller bins according to the second decimal digit and repeat.  Explain how continuing in this way produces digits $d_1, d_2, d_3, \ldots$ and nested closed intervals $I_1 \supseteq I_2 \supseteq I_3 \supseteq \cdots$, where $I_k$ has length $10^{-k}$ and contains $a_n$ for infinitely many $n$.
	\item Build a subsequence with $a_{n_k} \in I_k$ for each $k$, and prove that it converges.  What is the decimal expansion of its limit?
	\item Carry out the construction for the sequence
	\[ 0.4,\quad 0.41,\quad 0.414,\quad 0.4142,\quad \ldots \]
	whose $n$th term consists of the first $n$ decimal digits of $\sqrt{2} - 1$.  Every term of this sequence is rational.  Is the limit?
	\item Halves work and tenths work.  Would it also work to cut each interval into two \emph{unequal} pieces, say in the ratio $1:9$?  What is the one property of the cutting scheme that the proof really needs?
\end{enumerate}
\end{problem}

\begin{problem}[A hiker's steps]
\hint{For (b), group the steps.  The third and fourth together cover at least $\frac{1}{2}$, as do the fifth through eighth, and so on.  For (c), let her turn around whenever she has wandered too far.}
A hiker walks along the number line.  Her position after $n$ steps is $a_n$, so her $n$th step has length $|a_{n+1} - a_n|$.
\begin{enumerate}[(a)]
	\item Suppose that $|a_{n+1} - a_n| \leq 2^{-n}$ for every $n$.  Prove that $(a_n)$ is a Cauchy sequence, so that she is certainly headed \emph{somewhere}.
	\item Suppose instead that every step is forward and the $n$th step has length exactly $1/n$, so that $a_{n+1} = a_n + \frac{1}{n}$.  Show that she eventually passes every milestone, which is to say that the sequence $(a_n)$ is unbounded.
	\item Now let the $n$th step again have length exactly $1/n$, but allow her to choose its direction.  Describe a rule for choosing directions that keeps her inside the interval $[-1, 2]$ forever and yet leaves $(a_n)$ divergent.  Prove both claims.
	\item In (b) and (c) the steps shrink to zero, and still the sequence is not Cauchy.  Explain in a sentence or two, with reference to the quantifiers in the definition, why ``consecutive terms get close together'' is a weaker statement than ``the sequence is Cauchy.''
\end{enumerate}
\end{problem}

\begin{problem}[The sceptic's series]
\hint{For (c), expand $(\sqrt{x} - \sqrt{y})^2 \geq 0$.  For (e), apply the Cauchy Criterion for series to the block $a_{n+1} + a_{n+2} + \cdots + a_{2n}$.}
Suppose that $a_n \geq 0$ for all $n$ and that $\sum_{n=1}^{\infty} a_n$ converges.  A sceptical friend proposes the statements below.  Decide which are true, prove those, and give a counterexample to each of the others.
\begin{enumerate}[(a)]
	\item $\sum_{n=1}^{\infty} a_n^2$ converges.  (Does your answer change if the terms are allowed to be negative?)
	\item $\sum_{n=1}^{\infty} \sqrt{a_n}$ converges.
	\item $\sum_{n=1}^{\infty} \sqrt{a_n a_{n+1}}$ converges.
	\item $n a_n \to 0$.
	\item If in addition $(a_n)$ is decreasing, then $n a_n \to 0$.
\end{enumerate}
\end{problem}
